% !TeX program = xelatex
% UTF-8; compile with XeLaTeX from this directory or the repository root.
% Source: Google Drive / Notebook / bronchiolitis, in filename order.
% pic01.jpg: 20261016_053954660.jpg; Drive ID 1UeJwKWqkku5k7GiCjJ2n57mvu6QKfxZg
% pic02.jpg: 20261016_054002256.jpg; Drive ID 14YDal05WnoQMoKt9G8sW1Z1Xqcs2vu1P
% Dual-mode guard: compiles standalone with XeLaTeX, or is \input from the
% XePersian notebook (main.tex) where its local preamble below is skipped.
\newif\ifBronchiolitisStandalone
\ifdefined\settextfont
  \BronchiolitisStandalonefalse
\else
  \BronchiolitisStandalonetrue
\fi
\ifBronchiolitisStandalone
\documentclass[11pt,a4paper]{article}
\usepackage[margin=20mm]{geometry}
\usepackage{amsmath}
\usepackage{graphicx}
\usepackage{enumitem}
\usepackage{xepersian}
\settextfont{Amiri}
\setlatintextfont{DejaVu Serif}
\setlength{\parindent}{0pt}
\setlength{\parskip}{3pt}
\setlength{\emergencystretch}{3em}
\setlist[itemize]{itemsep=2pt,topsep=2pt}
\newcommand{\bmix}[1]{\lr{#1}}
\newcommand{\bheader}[1]{%
  \section*{تصویر #1}
  \textbf{موضوع:} ۲۰۹ ـ برونشیولیت\par
  \begingroup\small
  سربرگ چاپی: \bmix{Subject:}\quad\bmix{Year:}\quad\bmix{Month:}\quad\bmix{Date:}
  (خانه‌های تاریخ خالی‌اند). شماره‌های چاپی سطرها: ۱ تا ۹.
  \endgroup\par}

\begin{document}
\fi

% Helper macros used inside the body must exist in notebook mode too.
\providecommand{\bmix}[1]{\lr{#1}}
\providecommand{\bheader}[1]{%
  \section*{تصویر #1}
  \textbf{موضوع:} ۲۰۹ ـ برونشیولیت\par
  \begingroup\small
  سربرگ چاپی: \bmix{Subject:}\quad\bmix{Year:}\quad\bmix{Month:}\quad\bmix{Date:}
  (خانه‌های تاریخ خالی‌اند). شماره‌های چاپی سطرها: ۱ تا ۹.
  \endgroup\par}

\begin{center}
  {\Large\bfseries ۲۰۹ ـ برونشیولیت}\par
  \bmix{Bronchiolitis}
\end{center}

\bheader{۱}
\subsection*{تعریف، علت و انتقال}
\begin{itemize}
  \item برونشیولیت: درگیری برونشیول‌های کوچک، افزایش تولید مخاط، برونکواسپاسم گاهگاهی؛ می‌تواند انسداد راه هوایی بدهد.
  \item معمولاً دنبال عفونت ویروسی دستگاه تنفسی فوقانی؛ شایع‌ترین: \bmix{RSV}.
  \item شایع‌ترین سن: شیرخواران و کودکان کم‌سن؛ در شیرخواران تهدیدکنندهٔ حیات.
  \item عمده‌ترین راه انتقال: دست؛ دومین راه: آئروسل‌های تنفسی.
\end{itemize}

\subsection*{شیوع و گروه‌های پرخطر}
\begin{itemize}
  \item برونشیولیت شایع‌ترین علت بستری شیرخواران است؛ عمدتاً طی ۲ سال اول زندگی؛ پیک: ۲–۶ ماهگی.
  \item در شیرخواران سالم، درمان می‌تواند سرپایی باشد.
  \item گروه‌های پرخطر:
  \begin{itemize}
    \item شیرخوار نارس (۲۹ هفته).
    \item بیماری مزمن ریوی پره‌مچوریتی.
    \item بیماری قلبی مادرزادی با همودینامیک فعال.
    \item ضعف نوروماسکولار.
    \item ضعف ایمنی.
  \end{itemize}
  \item فصل شایع: اواخر زمستان.
\end{itemize}

\subsection*{برونشیولیت \bmix{RSV}}
\begin{itemize}
  \item دورهٔ انکوباسیون: ۴–۶ روز.
  \item به صورت یک بیماری تنفسی پیش‌رونده.
  \item سرفه + آبریزش بینی.
  \item ظرف ۳–۷ روز: تنفس \bmix{noisy}، ویز قابل شنیدن.
  \item معمولاً تب \bmix{low-grade} + تحریک‌پذیری.
  \item الگوی تنفس: بازدم طولانی، رتراکشن، نازال فلرینگ؛ در معاینهٔ ریه: ویز و کراکل منتشر.
  \item در موارد شدید: سیانوز و گرانتینگ.
\end{itemize}
\begingroup\small شمارهٔ دست‌نویس پایین صفحه: ۱.\endgroup

\clearpage
\bheader{۲}
\begin{itemize}
  \item در نوزاد و شیرخوار می‌تواند با آپنه شروع شود.
\end{itemize}

\subsection*{تشخیص}
\begin{itemize}
  \item برای تشخیص به آزمایشگاه نیاز نیست!
  \item \bmix{CXR} همیشه نیاز نیست.
  \item پرهوایی:
  \begin{itemize}
    \item افزایش لوسنسی ریه.
    \item تخت شدن دیافراگم.
  \end{itemize}
  \item اگر انفیلتره داشته: پنومونی ویروسی؛ آتلکتازی لوکالیزه.
\end{itemize}

\subsection*{\bmix{DDX}}
بیماری‌هایی که ویزینگ داشته:
\begin{itemize}
  \item آسم: سن ابتلا، تب، سابقهٔ خانوادگی آسم.
  \item جسم خارجی.
  \item \bmix{CF}.
  \item بیماری مزمن ریه ناشی از پره‌مچوریتی.
  \item \bmix{GERD}.
  \item پنومونی ویروسی یا باکتریال.
\end{itemize}

\subsection*{درمان}
\begin{itemize}
  \item حمایتی؛ ظرف ۲–۵ روز رو به بهبودی می‌روند.
  \item کنترل تب، هیدریشن، ساکشن مجاری هوایی فوقانی، \bmix{\(O_2\)} در صورت نیاز.
  \item اندیکاسیون بستری: دیسترس تنفسی، هایپوکسمی، آپنه، عدم تحمل \bmix{PO}.
  \item برونکودیلاتور، \bmix{AB}، کورتیکواستروئید، سالین هایپرتونیک، فیزیوتراپی تنفسی: کارایی ندارند.
  \item تاکی‌پنه و هایپوکسی: حرکت به سمت نارسایی تنفسی و نیاز به ونتیلاسیون مکانیکی.
\end{itemize}

\subsection*{پیشگیری}
\begin{itemize}
  \item \bmix{Palivizumab} ماهانه برای شیرخواران پرمعرض خطر؛ از یک ماه قبل از زمستان.
\end{itemize}
\begingroup\small شمارهٔ دست‌نویس پایین صفحه: ۲.
در حاشیهٔ چپ، کنار سطر «کارایی ندارند»، قطعه‌متن «…نتای» از برگهٔ زیرین دیده می‌شود؛ ادامهٔ آن پشت برگه پنهان است.
\endgroup

\clearpage
\section*{تصاویر اصلی}
\IfFileExists{images/pic01.jpg}{\def\bimagebase{images/}}{%
  \def\bimagebase{pediatrics/ped-pulmonology/bronchiolitis/images/}}
\begin{center}
  \textbf{تصویر ۱}\par
  \includegraphics[width=\linewidth,height=0.88\textheight,keepaspectratio]{\bimagebase pic01.jpg}
\end{center}
\clearpage
\begin{center}
  \textbf{تصویر ۲}\par
  \includegraphics[width=\linewidth,height=0.88\textheight,keepaspectratio]{\bimagebase pic02.jpg}
\end{center}
\ifBronchiolitisStandalone
\end{document}
\fi
